% ECONOMICS_PROBLEM: {"id": "D63-653", "title": "Complexity of verifying epistemic EFX for additive goods", "jel_code": "D63", "source_id": "OP-0193"}
% FIRST_POSTED: 2026-10-10
% ATTEMPT_STATUS: solved
\documentclass[11pt]{article}
\usepackage[margin=0.85in]{geometry}
\usepackage{amsmath,amssymb,amsthm}
\usepackage[T1]{fontenc}
\usepackage[hidelinks]{hyperref}
\hypersetup{pdftitle={Recognizing epistemic EFX allocations of goods},pdfauthor={Ji Ho Bae}}
\newtheorem{theorem}{Theorem}
\newcommand{\NP}{\mathsf{NP}}
\newcommand{\EEFX}{\mathrm{EEFX}}
\setlength{\parindent}{0pt}
\setlength{\parskip}{5pt}
% AI assistance: OpenAI Codex with automated collaboration.
% The mathematical draft was generated with AI assistance; no human proof edits were recorded.
% No priority claim is made beyond the cited sources.
\begin{document}
\begin{center}
{\Large\bfseries Recognizing epistemic EFX allocations of goods}\\[5pt]
{\small A short proof for catalogue problem D63-653}\\[3pt]
{\small Ji Ho Bae\quad---\quad10 October 2026}
\end{center}

An allocation $A=(A_1,\ldots,A_n)$ is an ordered partition of a finite
set of goods $M$. Agent $i$ has an additive valuation
$v_i(S)=\sum_{g\in S}v_i(g)$, explicitly encoded in binary.
Following Caragiannis et al.~\cite[Definitions 1--2]{fairness},
$A$ is \emph{epistemic EFX} (EEFX) if, for each $i$, there is an allocation
$Y^{(i)}$ with $Y^{(i)}_i=A_i$ such that
\[
 v_i(A_i)\geq v_i\bigl(Y^{(i)}_j\setminus\{g\}\bigr)
 \qquad\text{for every }j\ne i\text{ and }g\in Y^{(i)}_j
 \text{ with }v_i(g)>0.
\]
Different agents may use different witnesses. The final paragraph of
\cite[Section 4]{fairness} asks for the complexity of recognizing such an
allocation of goods. Its Section 6 treats chores separately.

\begin{theorem}
Deciding whether a supplied allocation is EEFX is NP-complete even for
exactly three agents with strictly positive integer item values, nonempty
supplied bundles, and a singleton supplied bundle for agent 1.
\end{theorem}

\begin{proof}
\textbf{Membership in $\NP$.}
A certificate supplies one witness allocation for each agent.
With $t$ goods it contains $nt$ recipient labels, each using
$O(\log n)$ bits. Check completeness, the fixed own bundle, and all EFX
inequalities. There are polynomially many checks, and binary integer
sums and comparisons take polynomial time. For nonnegative rational
values in the unrestricted problem, a common denominator has polynomial
bit length, so the same verifier remains polynomial.

\textbf{Reduction.}
Use the NP-complete problem \textsc{Partition}~\cite{karp}: given positive
integers $a_1,\ldots,a_m$ with sum $2B$, ask whether some subset sums to
$B$. This restricted source remains NP-complete. Indeed, a signed
equal-partition instance is unchanged by taking absolute values and
discarding zero entries; doubling all remaining entries makes the total
even. An all-zero input can be replaced by $(1,1)$.
Thus $m,B\geq1$.

Put $H=B+1$ and $V=mH+B$. Introduce $m$ regular goods
$g_1,\ldots,g_m$, $m+2$ dummy goods $z_1,\ldots,z_{m+2}$, and a good $d$.
Agent 1's item values are
\[
 v_1(g_\ell)=H+a_\ell,\qquad v_1(z_r)=H,\qquad v_1(d)=V.
\]
Supply the allocation
\[
 A_1=\{d\},\qquad
 A_2=\{g_1,\ldots,g_m,z_2,\ldots,z_{m+2}\},\qquad
 A_3=\{z_1\}.
\]
For each $j=2,3$, let $L=2m+3$ and set $v_j(g)=L$ if $g\in A_j$,
and $v_j(g)=1$ otherwise. Agent $j$ values its own bundle at least $L$
and its entire complement at most $L-1$. Therefore $A$ itself is an
EEFX witness for each of these two agents. Every item value is positive.

\textbf{Correctness.}
For positive additive values, EFX against a nonempty bundle $C$ is
equivalent to
\[
 v_1(C)-\min_{g\in C}v_1(g)\leq V,
 \tag{1}\label{eq:residual}
\]
because deleting a minimum-valued good leaves the largest residual.
An agent-1 witness partitions the $2m+2$ outside goods into two bundles.
A bundle with at least $m+2$ goods leaves at least $m+1$ goods after
any one deletion, giving a residual of at least
$(m+1)H>mH+B=V$. Hence both bundles have at most $m+1$ goods, and
counting forces each to have exactly $m+1$.

There are only $m$ regular goods, so each bundle contains a dummy and
has minimum value $H$. If $S_j$ indexes its regular goods, its residual is
\[
 mH+\sum_{\ell\in S_j}a_\ell.
\]
Condition \eqref{eq:residual} bounds both regular sums by $B$.
Together they sum to $2B$, so both equal $B$, giving a partition solution.

Conversely, divide the regular goods according to a partition solution,
and pad each side to $m+1$ goods with dummies. Each side needs at least
one dummy; the two sides need exactly $m+2$ in total. Their minimum
value is $H$, and both residuals equal $mH+B=V$. They give a witness
for agent 1, and the auxiliary witnesses already established show that
$A$ is EEFX.

\textbf{Size.}
The construction has $2m+3$ goods and three valuation rows.
Since $a_\ell\leq2B$, every value uses
$O(\log(m+1)+\log(B+1))$ bits. Constructing the explicit allocation and
valuation matrix therefore takes polynomial time. This proves the claim.
\end{proof}

\textbf{Formalization scope.}
The Lean companion verifies the original three-agent EEFX reduction
equivalence in \texttt{EEFX.reduction\_correct\_iff}, together with
strict positivity, the singleton own bundle, nonempty supplied bundles,
and the count $2m+3$ of goods. The complexity classes, the polynomial-time
encoding and verifier analysis, and the classical NP-completeness of
\textsc{Partition} are supplied by this hand proof and its cited source;
they are not claimed as Lean theorems.
% Public pinned source and reproduction links are appended in the catalogue submission.
\paragraph{Lean source and reproduction.}
\href{https://github.com/jbaelaw/eefx-recognition-lean/blob/4fa24ef0066295da70ef23370c9aedf87bd3114a/paper/ShortProof.pdf}{Proof PDF, Ji Ho Bae}.
The \href{https://github.com/jbaelaw/eefx-recognition-lean/tree/4fa24ef0066295da70ef23370c9aedf87bd3114a}{complete pinned Lean project}
contains \href{https://github.com/jbaelaw/eefx-recognition-lean/blob/4fa24ef0066295da70ef23370c9aedf87bd3114a/EEFX/Reduction.lean}{the main reduction theorem}
and the \href{https://github.com/jbaelaw/eefx-recognition-lean/blob/4fa24ef0066295da70ef23370c9aedf87bd3114a/original/initial-research-note.tex}{complete original generated note}.
Reproduction uses Lean 4.35.0-rc3 and
\href{https://github.com/leanprover-community/mathlib4/tree/b84a70d6a5ed793cc46184160f4d4188d8c825fb}{mathlib commit b84a70d6},
with all dependencies locked in the project manifest.
The remote commands \texttt{lake build} and \texttt{lake env lean Verify.lean}
both exited successfully, without warnings. There are no
\texttt{sorry}/\texttt{admit} placeholders or added axioms;
the main theorem uses only \texttt{propext}, \texttt{Classical.choice},
and \texttt{Quot.sound}.
\href{https://github.com/jbaelaw/eefx-recognition-lean/blob/4fa24ef0066295da70ef23370c9aedf87bd3114a/verification/LeanVerification.json}{The checking record}
identifies the exact scope and source hashes. Attribution is GPT Codex;
the exact serving model revision was not exposed. This is a complete
proof claim awaiting expert review, and is not a verification claim
by the catalogue maintainers.

\section*{Completion Estimate}
100\%
For the recognition-complexity question, as an unreviewed proof claim.
This estimate does not assert a Lean formalization of complexity classes.

\begin{thebibliography}{2}
\small
\bibitem{fairness}
I. Caragiannis, J. Garg, N. Rathi, E. Sharma, and G. Varricchio.
\emph{New Fairness Concepts for Allocating Indivisible Items}.
arXiv:2206.01710v3, 14 January 2025.
Definitions 1--2; Section 4, final paragraph.
\url{https://arxiv.org/html/2206.01710v3}.

\bibitem{karp}
R. M. Karp. \emph{Reducibility Among Combinatorial Problems}.
In \emph{Complexity of Computer Computations}, Plenum Press, 1972,
pp.~85--103. Section 4: Main Theorem (p.~94), problem 20 (p.~97),
and the Knapsack-to-Partition reduction (p.~100).
\href{https://doi.org/10.1007/978-1-4684-2001-2_9}{doi:10.1007/978-1-4684-2001-2\_9}.
\href{https://www.cs.umd.edu/~gasarch/BLOGPAPERS/Karp.pdf}{Original-page scan in the 2010 reprint}.
\end{thebibliography}
\end{document}
